% ============================================================
% PREAMBLE.TEX — Cấu hình packages và thiết lập chung
% Đồ án 1: Xây dựng và tối ưu mô hình YOLO26 cho thiết bị biên
% ============================================================

% --- Encoding & Ngôn ngữ ---
\usepackage[utf8]{inputenc}
\usepackage[T5]{fontenc}           % Font encoding cho tiếng Việt
\usepackage[vietnamese]{babel}     % Hỗ trợ ngôn ngữ Việt
\usepackage{csquotes}              % Load csquotes theo khuyến nghị của biblatex khi dùng babel

% --- Font & Math: Times New Roman & Mono ---
\usepackage{amsmath}               % Load amsmath trước các gói math khác
\usepackage{amssymb}               % Load amssymb trước newtxmath để tránh lỗi trùng lặp \Bbbk
\usepackage{tgtermes}              % Text font: TeX Gyre Termes (Times-compatible, hỗ trợ đầy đủ tiếng Việt không lỗi/cảnh báo)
\usepackage{newtxmath}             % Math font: Times-compatible
\usepackage[scaled=0.90]{zi4}      % Monospace font: Inconsolata (hỗ trợ đầy đủ các ký tự, tối ưu cho listings)

% --- TikZ & PGFPlots (Đồ họa nội sinh & Biểu đồ thống nhất) ---
\usepackage{tikz}
\usepackage{pgfplots}
\pgfplotsset{compat=1.18}
\usetikzlibrary{shapes.geometric, arrows.meta, positioning, calc, fit, backgrounds}
\tikzset{
    % Style cho sơ đồ khối
    block/.style={rectangle, draw=blue!70!black, fill=blue!5, thick, rounded corners, minimum width=2.2cm, minimum height=0.9cm, align=center, font=\small},
    orangeblock/.style={rectangle, draw=orange!80!black, fill=orange!5, thick, rounded corners, minimum width=2.2cm, minimum height=0.9cm, align=center, font=\small},
    greenblock/.style={rectangle, draw=green!60!black, fill=green!5, thick, rounded corners, minimum width=2.2cm, minimum height=0.9cm, align=center, font=\small},
    purpleblock/.style={rectangle, draw=purple!70!black, fill=purple!5, thick, rounded corners, minimum width=2.2cm, minimum height=0.9cm, align=center, font=\small},
    grayblock/.style={rectangle, draw=gray!70!black, fill=gray!10, thick, rounded corners, minimum width=2.2cm, minimum height=0.9cm, align=center, font=\small},
    line/.style={draw, -{Stealth[scale=1.0]}, thick},
    dashedline/.style={draw, -{Stealth[scale=1.0]}, thick, dashed},
    sum/.style={circle, draw=blue!70!black, fill=blue!5, thick, minimum size=0.5cm, inner sep=0pt, node contents={+}},
    splitnode/.style={circle, fill=black, inner sep=1.2pt}
}

% --- Trang & Layout ---
\usepackage[
    a4paper,
    left=3cm,
    right=2cm,
    top=2cm,
    bottom=2cm,
    headheight=20pt              % Tránh cảnh báo \headheight quá nhỏ từ fancyhdr
]{geometry}
\usepackage{setspace}
\onehalfspacing                    % Giãn dòng 1.5

% --- Tiêu đề chương/section ---
\usepackage{titlesec}
\titleformat{\chapter}[display]
    {\normalfont\Large\bfseries\centering}
    {\MakeUppercase{\chaptertitlename}\ \thechapter}{12pt}{\LARGE}
\titlespacing*{\chapter}{0pt}{-20pt}{20pt}

% --- Header / Footer ---
\usepackage{fancyhdr}
\pagestyle{fancy}
\fancyhf{}
\fancyhead[L]{\small\leftmark}
\fancyhead[R]{\small Đồ án 1 — Tối ưu YOLO26 cho thiết bị biên}
\fancyfoot[C]{\thepage}
\renewcommand{\headrulewidth}{0.4pt}

% --- Hình ảnh ---
\usepackage{graphicx}
\graphicspath{{figures/}}
\usepackage{float}                 % [H] placement
\usepackage{subcaption}            % subfigure

% --- Bảng ---
\usepackage{booktabs}              % \toprule, \midrule, \bottomrule
\usepackage{tabularx}
\usepackage{multirow}
\usepackage{array}
\usepackage{longtable}

% --- Toán ---
\usepackage{bm}                    % Bold math (amsmath và amssymb đã được load ở đầu file)
\usepackage{mathtools}

% --- Algorithm (dùng cho pseudo-code TAL trong Chương 5) ---
\usepackage{algorithm}
\usepackage{algpseudocode}
\floatname{algorithm}{Thuật toán}

% --- Code listing ---
\usepackage{listings}
\usepackage{xcolor}

\definecolor{codebg}{HTML}{F5F5F5}
\definecolor{codegreen}{HTML}{28A745}
\definecolor{codegray}{HTML}{6A737D}
\definecolor{codepurple}{HTML}{6F42C1}

\lstdefinestyle{pythonstyle}{
    language=Python,
    backgroundcolor=\color{codebg},
    basicstyle=\small\ttfamily,
    keywordstyle=\color{codepurple}\bfseries,
    stringstyle=\color{codegreen},
    commentstyle=\color{codegray}\itshape,
    numberstyle=\tiny\color{codegray},
    numbers=left,
    numbersep=8pt,
    breaklines=true,
    frame=single,
    rulecolor=\color{codegray!30},
    tabsize=4,
    showstringspaces=false,
    captionpos=b,
    morekeywords={self,True,False,None},
}
\lstset{style=pythonstyle}

% --- Hyperlinks ---
\usepackage[
    colorlinks=true,
    linkcolor=black,
    citecolor=blue!70!black,
    urlcolor=blue!70!black
]{hyperref}

% --- Tài liệu tham khảo ---
\usepackage[
    backend=biber,
    style=ieee,
    sorting=none,
    language=english             % Tránh lỗi/cảnh báo tiếng Việt không được biblatex hỗ trợ dịch nhãn
]{biblatex}
\addbibresource{bibliography.bib}

% --- Misc ---
\usepackage{enumitem}              % Tùy chỉnh list
\usepackage{caption}               % Tùy chỉnh caption
\captionsetup{
    font=small,
    labelfont=bf,
    skip=8pt
}
\usepackage{pdfpages}              % Chèn trang PDF
\usepackage{appendix}
\usepackage{framed}

% --- Định nghĩa hộp concept (lightweight, no tcolorbox) ---
\newenvironment{conceptbox}[1][]{%
    \begin{framed}\noindent\textbf{#1}\par\smallskip\noindent\ignorespaces
}{%
    \end{framed}
}

\newenvironment{notebox}[1][]{%
    \begin{framed}\noindent\textbf{#1}\par\smallskip\noindent\ignorespaces
}{%
    \end{framed}
}

% --- Custom commands ---
\usepackage{xspace}                % Load xspace trước khi định nghĩa lệnh dùng nó
\newcommand{\yolo}{YOLO26\xspace}
\newcommand{\rpi}{Raspberry Pi 5\xspace}
\newcommand{\figref}[1]{Hình~\ref{#1}}
\newcommand{\tabref}[1]{Bảng~\ref{#1}}
\newcommand{\eqnref}[1]{Công thức~(\ref{#1})}
\newcommand{\chapref}[1]{Chương~\ref{#1}}
\newcommand{\secref}[1]{Mục~\ref{#1}}
\newcommand{\algoref}[1]{Thuật toán~\ref{#1}}

